\section*{अध्याय \ref{ch:b3:extension-strike-slip} --- प्रसारी और स्ट्राइक-स्लिप विवर्तनिकी}

\begin{solution}{exo:b3:extension-strike-slip:1}
संस्तरों की नति $50 - 35 = \ang{15}$ है, भ्रंशों से दूर की ओर; $\beta =
\sin\ang{50}/\sin\ang{35} = 0.766/0.574 = 1.34$।
\end{solution}

\begin{solution}{exo:b3:extension-strike-slip:2}
दाएँ-पार्श्वीय, दाईं ओर सोपान: द्रोणी। दाएँ-पार्श्वीय, बाईं ओर सोपान: कटक। बाएँ-पार्श्वीय
बाईं ओर सोपान: द्रोणी। बाएँ-पार्श्वीय, दाईं ओर सोपान: कटक।
\end{solution}

\begin{solution}{exo:b3:extension-strike-slip:3}
$105/20 = \qty{5.25}{km/Myr} = \qty{5.25}{mm/yr}$ है, जो
\qty{6}{mm/yr} के निकट है: \qty{20}{Myr} की अवधि में दर का स्थिर रहना आवश्यक नहीं।
\end{solution}

\begin{solution}{exo:b3:extension-strike-slip:4}
\qty{200}{\percent} का अर्थ है $\Delta L = 2L_0$: $L_0 = \qty{117.5}{km}$ और $L
= \qty{352.5}{km}$। उत्तर में $\Delta L = 0.5L_0$: $L_0 = \qty{470}{km}$
और $L = \qty{705}{km}$।
\end{solution}

\begin{solution}{exo:b3:extension-strike-slip:5}
$563/20 = 28$ से $563/15 = \qty{38}{mm/yr}$ तक। यह विस्थापन का न्यूनतम मान है, और
भ्रंश क्षेत्र के आर-पार सर्वेक्षण से मापा गया सरकाव
प्रणाली के पड़ोसी भ्रंशों की गति और भूकंपों के बीच संचित प्रत्यास्थ
विकृति को भी सम्मिलित करता है।
\end{solution}

\begin{solution}{exo:b3:extension-strike-slip:6}
$t = 90 - \text{नति}$। \ang{60} पर: $(1 + 0.6\times1.732)/(1 - 0.6\times
0.577) = 2.039/0.654 = 3.12$; \ang{50} पर: $1.715/0.497 = 3.45$;
\ang{40} पर: $1.504/0.285 = 5.28$। लॉक-अप $\arctan0.6 = \ang{31.0}$ पर होता है।
\end{solution}

\begin{solution}{exo:b3:extension-strike-slip:7}
$\arctan0.4 = \ang{21.8}$ और $\arctan0.2 = \ang{11.3}$। \ang{15} नति वाला भ्रंश
केवल तभी फिसलता है जब $\mu < \tan\ang{15} = 0.27$।
\end{solution}

\begin{solution}{exo:b3:extension-strike-slip:8}
$2.5\times4.0 = \qty{10.0}{km^2}$ और $2.5\times(4.0 - 1.5) =
\qty{6.25}{km^2}$।
\end{solution}

\begin{solution}{exo:b3:extension-strike-slip:9}
$132/1.19 = \qty{110.9}{km}$ है; \qty{21.1}{km} का खिंचाव हुआ,
जिसमें $21.1/6 = \qty{3.5}{Myr}$ लगे।
\end{solution}

\begin{solution}{exo:b3:extension-strike-slip:10}
$3.0/\sin\ang{60} = \qty{3.464}{km}$ पहले और $3.0/\sin\ang{45} =
\qty{4.243}{km}$ बाद में: प्रति खंड \qty{0.779}{km} ($\beta = 1.22$)। संस्तरों की
नति $60 - 45 = \ang{15}$ है। प्रत्येक खंड एक इकाई के रूप में घूमता है: वह घूर्णन जो
पूर्व की ओर झुकने वाले, मान लीजिए, भ्रंशों की नति घटाता है, वह खंड की एक क्षैतिज रेखा को
इस प्रकार घुमाता है कि वह पूर्व की ओर ऊपर उठती है, और संस्तर पश्चिम की ओर नत होते हैं,
जो भ्रंशों की दिशा के विपरीत है।
\end{solution}

\begin{solution}{exo:b3:extension-strike-slip:11}
बाएँ-पार्श्वीय पूर्व--पश्चिम क्षेत्र में $\sigma_1$ का दिगंश \ang{045} है। R रीडेल \omterm{def:b3:deformation-mechanisms:shear-zone}{अपरूपण
क्षेत्र} से $\phi/2 = \ang{20}$ पर, $\sigma_1$ की ओर: स्ट्राइक \ang{070},
बाएँ-पार्श्वीय। R$'$ पर $90 - 20 = \ang{70}$: स्ट्राइक \ang{020}, दाएँ-पार्श्वीय। तनाव
विभंग $\sigma_1$ के समांतर: स्ट्राइक \ang{045}।
\end{solution}

\begin{solution}{exo:b3:extension-strike-slip:12}
$\mu = 0.6$ के लिए भ्रंश \ang{60.5} पर बनते हैं और \ang{31.0} पर लॉक होते हैं: $\beta =
\sin\ang{60.5}/\sin\ang{31.0} = 0.870/0.515 = 1.69$। $\mu = 0.85$ के लिए
\ang{65.2} से \ang{40.4} तक, $\beta = 0.908/0.648 = 1.40$। कम घर्षण वाले
भ्रंश लॉक होने से पहले अधिक घूमते हैं: दुर्बल चट्टानें नए
भ्रंश बनने से पहले अधिक खिंचती हैं।
\end{solution}

\begin{solution}{pb:b3:extension-strike-slip:1}
\textbf{1.} संस्तर पहले क्षैतिज थे और खंडों के साथ घूम गए हैं; भ्रंश
भी उतने ही कोण से घूमे हैं, अतः $\theta_0 = 42 + 18 = \ang{60}$,
$\phi = \ang{30}$ के लिए एंडरसन की नति।

\textbf{2.} $\beta = \sin\ang{60}/\sin\ang{42} = 0.8660/0.6691 = 1.294$।

\textbf{3.} $24.0/1.294 = \qty{18.5}{km}$; क्षेत्र
को \qty{5.5}{km} (\qty{5.46}{km}) से विस्तारित किया गया है।

\textbf{4.} $t = \ang{48}$: $(1 + 0.6\times0.900)/(1 - 0.6\times1.111) =
1.540/0.334 = 4.62$, जबकि \ang{60.5} पर नए भ्रंशों के लिए $3.12$ के मुकाबले
लगभग डेढ़ गुना अधिक।

\textbf{5.} $\arctan0.6 = \ang{31.0}$ पर; $\beta = \sin\ang{60}/\sin\ang{31}
= 1.68$ अधिक से अधिक। उससे पहले, लगभग \ang{60} पर नए भ्रंश खंडों को
काटते, और डोमिनो की दूसरी पीढ़ी उनका स्थान ले लेती।

\textbf{6.} ऊपरी भित्ति का प्रत्येक स्तंभ, क्षैतिज विस्थापन पर, अपने नीचे के भ्रंश के
गिराव जितना नीचे खिसकता है; भ्रंश सतह के निकट सबसे खड़ा है और उससे दूर
चपटा होता है, इसलिए भ्रंश के निकट के स्तंभ सबसे अधिक गिरते हैं, और
दूर वाले सबसे कम; संस्तर भ्रंश की ओर झुकते हैं।

\textbf{7.} $D = A/H = 16.0/3.2 = \qty{5.0}{km}$।

\textbf{8.} $D$, $A$ के समानुपाती है: $\qty{5.0}{km} \pm \qty{0.5}{km}$।

\textbf{9.} $3.2\times(5.0 - 2.0) = \qty{9.6}{km^2}$।

\textbf{10.} क्षैतिज विस्थापन, \qty{3.2}{km}।

\textbf{11.} भराव भ्रंश की ओर एक पच्चर के रूप में मोटा होता है, और इसकी
पुरानी परतें नई परतों की तुलना में अधिक झुकी हैं: प्रत्येक परत तब निक्षेपित हुई जब
रोलओवर और भ्रंश से सटा स्थान बढ़ता जा रहा था।

\textbf{12.} $z = P/(\rho g) = \num{5.5e8}/(2800\times9.8) = \qty{20.0}{km}$.

\textbf{13.} $20.0/(24 - 16) = \qty{2.5}{km/Myr} = \qty{2.5}{mm/yr}$.

\textbf{14.} $20.0/\sin\ang{30} = \qty{40}{km}$ का विस्थापन, और $40\cos\ang{30} =
\qty{34.6}{km}$ का क्षैतिज प्रसार।

\textbf{15.} नहीं: \ang{15} नति, \ang{31} की लॉक-अप नति से कम है।
या तो भ्रंश अधिक खड़ा रहते हुए सरका और फिर घुमाया गया (\omterm{def:b3:extension-strike-slip:core-complex}{लुढ़कता मोड़}), या
वह दुर्बल था, जिसमें उच्च दाब पर मृत्तिका या द्रवों से $\mu < 0.27$ था।

\textbf{16.} $45 - 15 = \ang{30}$ से: जैसे-जैसे ऊपरी प्लेट दूर खिसकी,
अभारित निचली भित्ति समस्थिति द्वारा ऊपर उठी और मुड़ी, और भ्रंश को
वलन-मोड़ के माध्यम से घुमा दिया।

\textbf{17.} \qty{34.6}{km} बनाम अर्ध-ग्राबेन का \qty{8.7}{km} (प्रश्न
23): चार गुना अधिक। \omterm{def:b3:extension-strike-slip:core-complex}{कायांतरी क्रोड संकुल} वहीं स्थित होते हैं जहाँ प्रसार
सबसे अधिक है, जैसे मध्य बेसिन ऐंड रेंज में।

\textbf{18.} बाईं ओर खिसकता बाएँ-पार्श्वीय भ्रंश: \omterm{def:b3:extension-strike-slip:bends}{विमोचक मोड़}, एक द्रोणी खुलती है।

\textbf{19.} $\qty{5}{mm/yr}\times\qty{1.2}{Myr} = \qty{6.0}{km}$.

\textbf{20.} $4.0\times6.0 = \qty{24}{km^2}$; $3.0 + 6.0 = \qty{9.0}{km}$ लंबी;
$9.0/4.0 = 2.25$.

\textbf{21.} $132/10 = 13$ से $132/7 = 19$: मृत सागर द्रोणी कहीं अधिक
लंबी है, उसका विस्थापन उसकी चौड़ाई का कई गुना है।

\textbf{22.} एक \omterm{def:b3:extension-strike-slip:bends}{अवरोधी मोड़}: शाखाओं के बीच ऊपर धकेला गया एक कटक; काट में,
एक धनात्मक \omterm{def:b3:extension-strike-slip:flower}{पुष्प संरचना}, जिसमें भ्रंश ऊपर की ओर विपाशित होकर बँटता है,
प्रतिलोम घटक वाली शाखाओं में, जैसी मृत सागर
रूपांतर पर 100 से \qty{300}{m} चौड़ी शाखा है।

\textbf{23.} $3.2 + 5.46 = \qty{8.66}{km}$: मेरो अर्ध-ग्राबेन
\textbf{\qty[text-series-to-math=true]{8.7}{km}} द्वारा विस्तारित हुआ है।
\end{solution}
